% ==========================================================================
%  Chapter 47 — Phase 2: Element Contract and Heterogeneous Containers
% ==========================================================================

\chapter{Phase 2: Element Contract and Heterogeneous Containers}
\label{ch:phase2-element-contract}

This chapter documents Phase~2 of the master plan
(Chapter~\ref{ch:convergence-master-plan}).  The original scope called
for introducing the \texttt{FiniteElement} C++20 concept, type-erased
element wrappers, and heterogeneous storage policies so that a single
\texttt{Model} can hold beams, shells, and continua simultaneously.
A deep codebase audit revealed that the element-contract and
heterogeneous-container infrastructure was \emph{already implemented}.
The single missing capability was the propagation of a
\texttt{revert\_material\_state()} operation through all layers of the
element architecture, mirroring the \texttt{commit\_material\_state()}
path documented in Phase~1 (Chapter~\ref{ch:phase1-constitutive-state}).
This chapter records the audit, the implementation of the revert chain,
and the verification test.

%% ------------------------------------------------------------------
\section{Phase 2 Scope Revision}
\label{sec:phase2-scope-revision}

Chapter~\ref{ch:convergence-master-plan} defined four Phase~2
deliverables:

\begin{enumerate}
  \item Define a C++20 concept \texttt{FiniteElement} with the
        assembly-level contract (\texttt{inject\_K},
        \texttt{compute\_internal\_forces},
        \texttt{inject\_tangent\_stiffness}, \texttt{commit},
        \texttt{revert}).
  \item Implement type-erased wrappers (\texttt{FEM\_Element},
        \texttt{StructuralElement}) via External Polymorphism.
  \item Implement \texttt{ElementPolicy}: homogeneous
        (\texttt{SingleElementPolicy<E>}) and heterogeneous
        (\texttt{MultiElementPolicy}).
  \item Integrate into \texttt{Model} with auto-detect, pre-built,
        and multi-material constructors.
\end{enumerate}

After a thorough audit of \texttt{src/elements/}, \texttt{src/model/},
and \texttt{src/analysis/}, deliverables~1--4 were found to be
\textbf{already implemented}, with one exception: the concept and all
wrappers enforced \texttt{commit\_material\_state(u)} but not
\texttt{revert\_material\_state()}.  The revert path is essential for
nonlinear solvers that require step bisection, line-search rejection,
or adaptive time stepping.

\begin{center}
\begin{tabular}{cll}
\toprule
\textbf{\#} & \textbf{Deliverable} & \textbf{Status before Phase 2} \\
\midrule
1 & \texttt{FiniteElement} concept &
    Implemented, missing \texttt{revert\_material\_state} \\
2 & Type-erased wrappers &
    Implemented, missing revert forwarding \\
3 & \texttt{ElementPolicy} storage &
    Fully implemented \\
4 & \texttt{Model} constructors &
    Fully implemented (3 constructors) \\
\bottomrule
\end{tabular}
\end{center}

Consequently, Phase~2 was re-scoped to:
\begin{enumerate}[label=(\alph*)]
  \item Comprehensive audit and documentation of the existing element
        architecture.
  \item Implementation of the full \texttt{revert\_material\_state}
        chain from the analysis boundary down to the constitutive
        storage layer.
  \item A verification test exercising every layer of the revert
        chain, including heterogeneous containers.
\end{enumerate}

%% ------------------------------------------------------------------
\section{The Element Contract}
\label{sec:finite-element-concept}

The assembly-level contract for every element type is expressed as a
C++20 concept in \texttt{FiniteElementConcept.hh}:

\begin{lstlisting}[caption={The \texttt{FiniteElement} concept},
                   label={lst:finite-element-concept}]
template <typename E>
concept FiniteElement =
  requires(E e, E const ce, Mat K, Vec u, Vec f) {
    { ce.num_nodes()              } -> std::convertible_to<std::size_t>;
    { ce.num_integration_points() } -> std::convertible_to<std::size_t>;
    { ce.sieve_id()               } -> std::convertible_to<PetscInt>;
    e.set_num_dof_in_nodes();
    e.inject_K(K);
    e.compute_internal_forces(u, f);
    e.inject_tangent_stiffness(u, K);
    e.commit_material_state(u);
    e.revert_material_state();                 // Phase 2 addition
  };
\end{lstlisting}

The key design decisions are:

\begin{itemize}
  \item \textbf{Commit takes displacement.}
        \texttt{commit\_material\_state(u)} receives the global solution
        vector because commit requires evaluation of the kinematic chain
        (shape gradients, deformation gradient, strains) before updating
        the constitutive state.
  \item \textbf{Revert takes nothing.}
        \texttt{revert\_material\_state()} is argument-free: it restores
        the last successfully committed state without needing kinematics,
        because all information resides in the committed storage.
  \item \textbf{Compile-time enforcement.}
        Any type that fails to provide either operation is rejected at
        compile time, preventing silent omission.
\end{itemize}

\subsection{Formal State Semantics}
\label{sec:commit-revert-semantics}

Let $\mathcal{S} = (\bsigma, \bC^{\text{ep}}, \bm{\alpha})$ denote
the full constitutive state at an integration point (stress, tangent,
internal variables).  The element maintains a committed snapshot
$\mathcal{S}^n$ and a trial state $\mathcal{S}^{\text{trial}}$ at each
material site.  The two operations satisfy:

\begin{align}
  \texttt{commit}(\mathbf{u}_{n+1}):
  \quad & \mathcal{S}^n \leftarrow
          \mathcal{S}^{\text{trial}}(\mathbf{u}_{n+1}),
  \label{eq:commit-op} \\[4pt]
  \texttt{revert}():
  \quad & \mathcal{S}^{\text{trial}} \leftarrow \mathcal{S}^n.
  \label{eq:revert-op}
\end{align}

Together they form an \emph{atomic checkpoint}:
\[
  \texttt{revert}(\texttt{commit}(\mathbf{u})) = \mathcal{S}^n,
  \qquad
  \texttt{commit}(\texttt{revert}()) = \mathcal{S}^n.
\]

This invariant is relied upon by Newton line-search strategies that
tentatively evaluate residuals without permanently altering the
history.

%% ------------------------------------------------------------------
\section{External Polymorphism: Type-Erased Wrappers}
\label{sec:type-erased-wrappers}

To hold elements of different concrete types (beams, shells, continua)
in a homogeneous container, the library uses the \emph{External
Polymorphism} pattern~\cite{Iglberger2022} implemented as a
Concept/Model/\texttt{pimpl} triple.

\subsection{FEM\_Element}
\label{sec:fem-element}

\texttt{FEM\_Element} is the assembly-level type-erased wrapper:

\begin{lstlisting}[caption={\texttt{FEM\_Element} type-erasure structure},
                   label={lst:fem-element}]
class FEM_Element {
  // --- Inner concept (virtual interface) ---
  struct Concept {
    virtual ~Concept() = default;
    virtual std::size_t num_nodes() const = 0;
    virtual std::size_t num_integration_points() const = 0;
    // ... inject_K, compute_internal_forces, etc.
    virtual void commit_material_state(Vec u) = 0;
    virtual void revert_material_state() = 0;  // Phase 2
  };

  // --- Model bridges to concrete element ---
  template <FiniteElement T>
  struct Model final : Concept {
    T element_;
    void commit_material_state(Vec u) override {
      element_.commit_material_state(u);
    }
    void revert_material_state() override {
      element_.revert_material_state();         // Phase 2
    }
    // ...
  };

  std::unique_ptr<Concept> pimpl_;

public:
  // --- Public API ---
  void commit_material_state(Vec u) { pimpl_->commit_material_state(u); }
  void revert_material_state()      { pimpl_->revert_material_state(); }
};
\end{lstlisting}

A \texttt{static\_assert} at file scope verifies that
\texttt{FEM\_Element} itself satisfies the \texttt{FiniteElement}
concept, closing the recursive contract.

\subsection{StructuralElement}
\label{sec:structural-element}

\texttt{StructuralElement} extends the same Concept/Model/\texttt{pimpl}
triple with structural-engineering introspection:

\begin{itemize}
  \item \texttt{concrete\_type()} returns a
        \texttt{std::type\_info const\&} for RTTI-based dispatch.
  \item \texttt{as<T>()} recovers the concrete element pointer for
        post-processing (section reconstruction, VTK export).
  \item \texttt{raw\_ptr()} exposes the void pointer for low-level
        operations.
\end{itemize}

The revert chain is identical to \texttt{FEM\_Element}: a virtual
\texttt{revert\_material\_state()} on the inner concept, forwarded
through the model to the concrete element.

%% ------------------------------------------------------------------
\section{Element Storage Policies}
\label{sec:element-policies}

\texttt{ElementPolicy.hh} defines two storage strategies constraining
how \texttt{Model} holds its element collection:

\begin{lstlisting}[caption={Element storage policies},
                   label={lst:element-policies}]
template <FiniteElement ElementT>
struct SingleElementPolicy {
  using element_type   = ElementT;
  using container_type = std::vector<ElementT>;
  static constexpr bool is_homogeneous = true;
};

struct MultiElementPolicy {
  using element_type   = FEM_Element;
  using container_type = std::vector<FEM_Element>;
  static constexpr bool is_homogeneous = false;
};
\end{lstlisting}

Both are constrained by the \texttt{ElementPolicyLike} concept:

\begin{lstlisting}
template <typename P>
concept ElementPolicyLike = requires {
  typename P::element_type;
  typename P::container_type;
  { P::is_homogeneous } -> std::convertible_to<bool>;
} && FiniteElement<typename P::element_type>;
\end{lstlisting}

\begin{remark}[Zero-overhead homogeneous path]
When the model contains a single element family (e.g.\ only
\texttt{ContinuumElement<Hex8,...>}), the \texttt{SingleElementPolicy}
stores concrete elements directly, avoiding virtual dispatch entirely.
The \texttt{MultiElementPolicy} pays only one pointer indirection per
element through the \texttt{pimpl} in \texttt{FEM\_Element}.
\end{remark}

%% ------------------------------------------------------------------
\section{Model Constructors}
\label{sec:model-constructors}

The \texttt{Model<Policy>} class supports three construction modes:

\begin{description}
  \item[Constructor 1 --- Auto-detect]
  Takes a \texttt{Domain} and a material--integrator pair.  Infers
  the element family from the mesh topology, constructs all elements
  with an identical constitutive model.

  \item[Constructor 2 --- Pre-built container]
  Takes a \texttt{Domain} and a ready-made \texttt{container\_type}.
  Useful when elements are constructed externally (e.g.\ from a script
  or a multi-physics coupler).

  \item[Constructor 3 --- Multi-material by physical group]
  Takes a \texttt{Domain} and a map from physical-group labels to
  material--integrator pairs.  Constructs heterogeneous elements
  (beam + shell + continuum) from a single mesh, with distinct
  constitutive models per group.
\end{description}

All three constructors produce models that support uniform iteration
over elements for assembly, commit, and revert.

%% ------------------------------------------------------------------
\section{The Revert Chain}
\label{sec:revert-chain}

\subsection{Architectural Layers}
\label{sec:revert-layers}

The revert operation propagates through six layers from the solver
boundary down to the constitutive storage:

\begin{center}
\begin{tikzpicture}[
  layer/.style={draw, rounded corners=3pt, minimum width=10cm,
                minimum height=0.75cm, font=\small\ttfamily},
  arr/.style={-{Stealth[length=5pt]}, thick},
  node distance=0.6cm
]
  \node[layer, fill=blue!8]  (L1) {Analysis / SNES boundary};
  \node[layer, fill=blue!12, below=of L1] (L2)
    {Element::revert\_material\_state()};
  \node[layer, fill=blue!16, below=of L2] (L3)
    {MaterialPoint::revert()  /  MaterialSection::revert()};
  \node[layer, fill=blue!20, below=of L3] (L4)
    {Material<P>::revert() \textrm{(type-erased)}};
  \node[layer, fill=blue!24, below=of L4] (L5)
    {ConstitutiveIntegrator::revert(site)};
  \node[layer, fill=blue!28, below=of L5] (L6)
    {TrialCommittedState::revert\_trial()};

  \draw[arr] (L1) -- (L2);
  \draw[arr] (L2) -- (L3);
  \draw[arr] (L3) -- (L4);
  \draw[arr] (L4) -- (L5);
  \draw[arr] (L5) -- (L6);

  \node[right=0.3cm of L1, font=\small\itshape, text=gray] {Phase 3};
  \node[right=0.3cm of L2, font=\small\itshape, text=gray] {Phase 2};
  \node[right=0.3cm of L3, font=\small\itshape, text=gray] {Phase 2};
  \node[right=0.3cm of L4, font=\small\itshape, text=gray] {Phase 2};
  \node[right=0.3cm of L5, font=\small\itshape, text=gray] {Phase 2};
  \node[right=0.3cm of L6, font=\small\itshape, text=gray] {Phase 1};
\end{tikzpicture}
\end{center}

Layers~2--5 were implemented in this phase.  Layer~6 existed from
Phase~1.  Layer~1 (the solver-level revert trigger from SNES convergence
monitoring or arc-length step rejection) belongs to Phase~3.

\subsection{Layer-by-Layer Trace}
\label{sec:revert-trace}

\paragraph{Layer 2: Element.}
Each concrete element iterates over its constitutive sites:
\begin{lstlisting}
// ContinuumElement
void revert_material_state() {
    for (auto& mp : material_points_) mp.revert();
}

// BeamElement, TimoshenkoBeamN, ShellElement
void revert_material_state() {
    for (auto& section : sections_) section.revert();
}
\end{lstlisting}

\paragraph{Layer 3: Constitutive site.}
Both \texttt{MaterialPoint} and \texttt{MaterialSection} delegate
directly:
\begin{lstlisting}
void revert() { material_.revert(); }
\end{lstlisting}

\paragraph{Layer 4: Type-erased material handle.}
\texttt{Material<P>} forwards through its \texttt{pimpl}:
\begin{lstlisting}
class Material {
  struct MaterialConcept {
    virtual void revert() = 0;       // pure virtual
  };
  template <typename Mat, typename Integ>
  struct OwningMaterialModel : MaterialConcept {
    void revert() override {
      integrator_.revert(material_);  // delegate to integrator
    }
  };
  void revert() { pimpl_->revert(); }
};
\end{lstlisting}

\paragraph{Layer 5: Constitutive integrator.}
The integrator performs a two-step revert:
\begin{enumerate}
  \item \textbf{Relation revert}: if the constitutive relation holds
        internal history (plasticity, damage), call
        \texttt{relation.revert(algo\_state)} to undo the provisional
        update.
  \item \textbf{State revert}: call
        \texttt{revert\_constitutive\_state(site)} which invokes
        \texttt{site.constitutive\_state().revert\_trial()}.
\end{enumerate}

\begin{lstlisting}[caption={Revert in the three integrators},
                   label={lst:integrator-revert}]
// PassThroughIntegrator (elastic)
void revert(auto& site) {
    // No-op: elastic materials have no history to revert.
}

// EmbeddedRelationIntegrator (inelastic)
void revert(auto& site) {
    if constexpr (requires { ... })
        site.constitutive_relation().revert(
            site.algorithmic_state());
    revert_constitutive_state(site);
}

// ContinuumLocalProblemIntegrator
void revert(auto& site) {
    if constexpr (requires { ... })
        constitutive_problem_.constitutive_relation().revert(
            site.algorithmic_state());
    revert_constitutive_state(site);
}
\end{lstlisting}

\paragraph{Layer 6: Storage.}
\texttt{TrialCommittedState<S>::revert\_trial()} copies the committed
snapshot back over the trial state, undoing any provisional update.
This layer was already present from Phase~1.

\subsection{Free Functions}

The constitutive integrator layer provides the revert free function in
the \texttt{constitutive\_integrators} namespace:

\begin{lstlisting}
template <typename ConstitutiveSiteT>
inline void revert_constitutive_state(ConstitutiveSiteT& site) {
    if constexpr (requires {
        site.constitutive_state().revert_trial(); })
        site.constitutive_state().revert_trial();
}
\end{lstlisting}

The \texttt{if constexpr} guard ensures elastic sites (with
committed-only storage) compile silently as no-ops.

%% ------------------------------------------------------------------
\section{Files Modified}
\label{sec:files-modified}

\begin{center}
\small
\begin{tabular}{lp{8.5cm}}
\toprule
\textbf{File} & \textbf{Change} \\
\midrule
\texttt{FiniteElementConcept.hh} &
  Added \texttt{e.revert\_material\_state()} to the concept
  requirements. \\
\texttt{FEM\_Element.hh} &
  Added \texttt{virtual void revert\_material\_state() = 0} to Concept;
  forwarding in Model; public API delegation. \\
\texttt{StructuralElement.hh} &
  Same pattern as \texttt{FEM\_Element}. \\
\texttt{ContinuumElement.hh} &
  Added \texttt{revert\_material\_state()} iterating over
  \texttt{material\_points\_}. \\
\texttt{BeamElement.hh} &
  Added \texttt{revert\_material\_state()} iterating over
  \texttt{sections\_}. \\
\texttt{TimoshenkoBeamN.hh} &
  Same pattern as \texttt{BeamElement}. \\
\texttt{ShellElement.hh} &
  Same pattern as \texttt{BeamElement}. \\
\texttt{MaterialPoint.hh} &
  Added \texttt{void revert()} forwarding to \texttt{Material<>}. \\
\texttt{MaterialSection.hh} &
  Added \texttt{void revert()} forwarding to \texttt{Material<>}. \\
\texttt{Material.hh} &
  Added \texttt{virtual void revert() = 0} to \texttt{MaterialConcept};
  implemented in \texttt{OwningMaterialModel},
  \texttt{NonOwningMaterialModel}, \texttt{MaterialRef},
  \texttt{Material<>}. \\
\texttt{ConstitutiveIntegrator.hh} &
  Added \texttt{revert\_constitutive\_state()} free function;
  added \texttt{revert()} to all three integrators. \\
\texttt{CMakeLists.txt} &
  Added \texttt{fall\_n\_heterogeneity\_test} target and CTest
  registration. \\
\bottomrule
\end{tabular}
\end{center}

%% ------------------------------------------------------------------
\section{Verification Test}
\label{sec:verification-test}

The file \texttt{tests/test\_element\_heterogeneity.cpp} contains
22~assertions organized in eight test groups:

\begin{center}
\small
\begin{tabular}{clp{7cm}}
\toprule
\textbf{Test} & \textbf{Name} & \textbf{What it validates} \\
\midrule
1 & Concept enforcement &
    \texttt{static\_assert} that a mock element satisfying all
    requirements passes \texttt{FiniteElement}; a mock missing
    \texttt{revert\_material\_state} is rejected. \\
2 & Wrapper forwarding &
    \texttt{FEM\_Element} and \texttt{StructuralElement} forward
    \texttt{revert\_material\_state()} to the wrapped concrete type. \\
3 & Heterogeneous container &
    \texttt{MultiElementPolicy} stores elements of different concrete
    types; uniform commit and revert iteration succeeds. \\
4 & Material type-erasure &
    \texttt{Material<>} revert through the erased boundary: elastic
    material fully restores tangent; plasticity material calls revert
    safely (no crash). \\
5 & Elastic safety &
    Elastic revert is a safe no-op: the tangent remains unchanged
    after \texttt{commit} followed by \texttt{revert}. \\
6 & ContinuumElement chain &
    Full commit/revert through a \texttt{ContinuumElement} with
    a PETSc mesh, verifying that \texttt{MaterialPoint::revert()}
    exercises the complete 6-layer chain. \\
7 & SNES incremental solve &
    Two-step nonlinear solve with \texttt{SNES}; after each
    converged step, \texttt{commit} then \texttt{revert} is called
    to verify the post-solve revert contract. \\
8 & Policy type traits &
    \texttt{static\_assert} coverage for
    \texttt{SingleElementPolicy}, \texttt{MultiElementPolicy},
    and the \texttt{ElementPolicyLike} concept. \\
\bottomrule
\end{tabular}
\end{center}

All 22 assertions pass.  The test executable is registered as CTest
target \texttt{element\_heterogeneity}.

%% ------------------------------------------------------------------
\section{Known Gaps}
\label{sec:phase2-gaps}

\begin{enumerate}
  \item \textbf{PlasticityRelation internal revert.}
        The \texttt{PlasticityRelation::revert(algo\_state)} method
        reverts the algorithmic state (accumulated plastic strain,
        back-stress), but the relation's \emph{own} internal members
        (e.g.\ the Newton iteration scratch) are not part of the
        \texttt{AlgorithmicState} object.  Completing this path is
        delegated to Phase~3 (Analysis Loop).

  \item \textbf{Analysis-level revert trigger.}
        The solver boundary (Layer~1) does not yet call
        \texttt{revert\_material\_state()} on convergence failure.
        Phase~3 will close this by adding step-rejection logic in the
        nonlinear analysis driver.

  \item \textbf{Dynamic history variables.}
        For time-integration schemes (Newmark, generalized-$\alpha$),
        velocities and accelerations at material sites must also be
        reverted.  This is delegated to Phase~5 (Dynamic Extension).
\end{enumerate}

%% ------------------------------------------------------------------
\section{Build and Test Results}
\label{sec:phase2-build-results}

After implementing the revert chain and all supporting changes:

\begin{center}
\begin{tabular}{lr}
\toprule
\textbf{Metric} & \textbf{Value} \\
\midrule
CMake targets compiled    & 27/27 (zero errors) \\
CTest results             & 41/43 pass \\
Pre-existing failures     & 2 (cantilever\_benchmark, snes\_gmsh\_pipeline) \\
New regressions           & 0 \\
Phase~2 verification test & 22/22 assertions pass \\
\bottomrule
\end{tabular}
\end{center}

%% ------------------------------------------------------------------
\section{Summary}
\label{sec:phase2-summary}

Phase~2 confirmed that the \texttt{fall\_n} element architecture already
provides a complete heterogeneous-container infrastructure: a C++20
\texttt{FiniteElement} concept, External Polymorphism wrappers
(\texttt{FEM\_Element}, \texttt{StructuralElement}), zero-overhead and
type-erased storage policies (\texttt{SingleElementPolicy},
\texttt{MultiElementPolicy}), and three \texttt{Model} construction
modes.  The only missing piece --- the \texttt{revert\_material\_state}
operation --- was implemented through all six architectural layers, from
the concrete element loops down to the
\texttt{TrialCommittedState::revert\_trial()} storage primitive.

Twenty-two verification assertions confirm the contract across concept
enforcement, type-erasure forwarding, heterogeneous iteration, and an
end-to-end SNES incremental solve.

Phase~3 will close the loop by wiring the analysis driver to call
\texttt{revert\_material\_state()} on step rejection, enabling robust
nonlinear continuation with bisection and adaptive time stepping.
